\documentclass[11pt,a4paper]{article}
\usepackage[margin=1in]{geometry}
\usepackage{hyperref}
\usepackage{graphicx}
\usepackage{enumitem}
\usepackage{booktabs}
\usepackage{siunitx}
\usepackage{longtable}
\usepackage{listings}
\usepackage{parskip}

\title{Piezo Experiment Software User Manual}
\author{Piezo Experiment Team}
\date{\today}

\begin{document}
\maketitle

\section{Overview}
This document explains the functionalities, main features, and usage of the Piezo Experiment software (Python/Tkinter UI). It also compares the current software to the legacy C\# implementation.

The software provides an end-to-end workflow for planning measurements, running the experiment (motor + Red Pitaya acquisition), monitoring live data, and exporting analysis outputs.

\section{System Requirements}
\begin{itemize}[leftmargin=*]
\item Operating system: Linux (current development) or Windows (packaged executable).
\item Python 3.x with required libraries when running from source.
\item Red Pitaya reachable over TCP/IP.
\item EPOS motor controller (optional for motor control modes).
\end{itemize}

\section{Quick Start}
\begin{enumerate}
\item Launch the UI: \texttt{python3 ui6.py} (or double-click the Windows \texttt{piezo-experiment.exe}).
\item Fill the experiment parameters in the \textbf{Run} tab (sample name, frequency and resistance ranges, repetitions, etc.).
\item Select the test mode (Plan-only, Pitaya-only, EPOS-only, or Full experiment).
\item Click \textbf{Run}. The software generates the plan, runs the measurement loop, and stores outputs under \texttt{generated\_configs/}.
\item Use the \textbf{Live acquisition} tab to monitor the latest CSV and view pressure/voltage p2p in real time.
\item Use the \textbf{Analysis} tab to load summary tables and export results.
\end{enumerate}

\section{User Interface Tour}
\subsection{Run Tab}
The Run tab defines the experiment parameters and starts the workflow.
\begin{itemize}[leftmargin=*]
\item Sample name and output folder parameters.
\item Frequency sweep (min/max, points, scale).
\item Resistance sweep (min/max, points, scale).
\item Periods and repetitions.
\item Pressure calibration (Vmin/Vmax, Pmin/Pmax).
\item 10x probe toggle.
\item Instrument input resistance (used for total resistance and power computation).
\item Pitaya host/port.
\item Mode selector: Plan-only, Pitaya-only, EPOS-only, Full experiment.
\item Safety controls: RPM limit and number of points.
\item DR08 prompt checkbox and digit display.
\end{itemize}

\subsection{Log Tab}
The Log tab shows the command output from the generator and executor, including errors and progress messages.

\subsection{Live Acquisition Tab}
The Live Acquisition tab plots the latest measurement CSV. It displays:
\begin{itemize}[leftmargin=*]
\item Current point, frequency, resistance, and repetition.
\item Live pressure and voltage traces.
\item Live p2p and mean statistics computed directly from the CSV.
\item Controls for stopping or continuing when prompted.
\end{itemize}

\subsection{Analysis Tab}
The Analysis tab loads and displays the summary table and allows exporting it. It also shows plots generated during data treatment.

\section{Workflow Details}
\subsection{Plan Generation}
The software generates an experiment plan using \texttt{gen\_config.py}, creating a run directory under \texttt{generated\_configs/}. It writes:
\begin{itemize}[leftmargin=*]
\item \texttt{experiment\_config.json}
\item \texttt{measurement\_plan.json}
\item \texttt{pitaya\_setup\_first\_point.txt}
\end{itemize}

\subsection{Execution}
The executor reads the plan and runs the experiment:
\begin{itemize}[leftmargin=*]
\item EPOS motor control (optional).
\item Red Pitaya acquisition (optional).
\item DR08 prompts when resistance changes (optional).
\item CSV output for each measurement under \texttt{run\_dir/csv/}.
\item A manifest with metadata (\texttt{run\_manifest.json}).
\end{itemize}

\subsection{Data Treatment and Summary}
The UI produces \texttt{summary\_table.csv} by reading each CSV and computing:
\begin{itemize}[leftmargin=*]
\item Pressure peak-to-peak.
\item Voltage peak-to-peak.
\item Mean power (from the CSV \texttt{power} column).
\end{itemize}
The plots \texttt{CombinedData\_PressureGraph.png}, \texttt{CombinedData\_VoltageGraph.png}, and \texttt{CombinedData\_PowerGraph.png} are generated for quick inspection.

\section{Outputs and File Structure}
Each run creates a directory like:
\begin{verbatim}
/generated_configs/YYYY_MM_DD_HHhMM/<sample_name>/
  experiment_config.json
  measurement_plan.json
  pitaya_setup_first_point.txt
  csv/
    SingleMeasurement_<freq>Hz_<R>Ohms_<rep>.csv
  summary_table.csv
  CombinedData_*.png
  run_manifest.json
\end{verbatim}

\section{Power Computation}
Power is computed per sample as:
\[ P = \frac{V^2}{R_{\text{total}}} \]
where:
\begin{itemize}[leftmargin=*]
\item $V$ is the voltage after calibration and probe scaling.
\item $R_{\text{total}}$ is the parallel combination of the load resistance and the instrument input resistance.
\end{itemize}
The instrument input resistance is configurable in the Run tab so users can update it when using a different multimeter or acquisition device.

\section{Troubleshooting}
\begin{itemize}[leftmargin=*]
\item \textbf{No CSV files appear:} verify Pitaya connection and that the selected mode includes acquisition.
\item \textbf{Live plot not updating:} ensure a CSV exists in the run directory and click Refresh.
\item \textbf{Motor not moving:} ensure EPOS is enabled in the selected mode and hardware is connected.
\item \textbf{Permission errors:} ensure the run directory is writable.
\end{itemize}

\section{Comparison with Legacy C\# Software}
The legacy application (C\# \texttt{Form1}) was a monolithic Windows Forms implementation. Key differences:
\begin{itemize}[leftmargin=*]
\item \textbf{Modularity:} the new software separates planning (\texttt{gen\_config.py}), execution (\texttt{executor.py}), and UI (\texttt{ui6.py}). The legacy code combined all responsibilities in a single form.
\item \textbf{Parameter flexibility:} the current UI exposes more parameters directly (e.g., signal multiplier, pressure calibration, instrument input resistance). In the legacy code many constants were fixed (e.g., Pitaya hostname, input impedance assumptions).
\item \textbf{Performance and robustness:} the new system isolates long-running acquisition in a worker thread and provides clear logging and retries; the legacy app ran the loop inside UI handlers and could become unresponsive.
\item \textbf{File outputs:} the new system standardizes run directories and writes JSON configs, measurement plans, per-point CSVs, and summary tables. The legacy application wrote CSVs and plots but did not emit structured configuration files.
\item \textbf{Safety controls:} explicit RPM caps and mode-specific safety checks are enforced in the new software.
\item \textbf{Cross-platform:} the new system can be packaged for Windows as a standalone executable; the legacy was tied to C\#/.NET Windows Forms.
\end{itemize}

\section{Appendix: Build (Windows Executable)}
To build a Windows executable from source, use PyInstaller on Windows:
\begin{verbatim}
python -m venv .venv
.venv\Scripts\activate
pip install -r requirements.txt
pip install pyinstaller
pyinstaller --onefile --noconsole --name piezo-experiment ui6.py
\end{verbatim}

\end{document}
